\section{Discussion}\label{sec:discussion}

\subsection{Hypotheses}
\textbf{H1 (sub-linear scaling, limited by physical cores): supported.} The best speedup on two physical cores
is 2.12, and almost all of it comes from the second core: two threads on two cores already reach 1.98, while
the second hyperthread of each core adds only 7\%. Hyper-Threading~\cite{marr2002} helps little because the
dominant Gaussian passes keep both hardware threads of a core waiting on the same load and floating-point
units. With close binding the curve shows the topology directly: two threads on one core gain 7--8\%, the same
Hyper-Threading gain seen at $p=4$. The Karp--Flatt fraction rises from 0.31 at $p=4$ to 0.42 at $p=8$, which
indicates growing overhead from oversubscription rather than a fixed serial part~\cite{karp1990}. The high
value at $p=2$ (0.86) reflects the close placement, not serial code: with spread placement the efficiency at
$p=2$ is 0.99.

\textbf{H2 (speedup stable or growing with image size): supported.} The 4-thread speedup stays between 2.13 and
2.31 from 0.33 to 21.1 megapixels and is slightly higher above native size, where each parallel region has more
work to amortise its fork and join. The serial cost per megapixel rises by 24\% at 21.1 megapixels. At that size
the six float images used per scale (24 bytes per pixel, about 500\,MB) far exceed the 55\,MiB L3, so the
column passes stream from DRAM; the parallel version suffers the same effect, which is why the speedup does not
change~\cite{williams2009}.

\textbf{H3 (best decomposition differs by stage): supported.} In CLAHE, one thread per tile needs no
synchronisation, and the 64 tiles cost almost the same, so dynamic scheduling neither helps nor hurts. The
atomic variant is 53\% slower even on one thread, because every increment becomes a locked read-modify-write,
and gets worse with more threads. In the Frangi stage the cost of a scale grows with its kernel length
($2\lceil3\sigma\rceil+1$ = 7, 13, 19 and 25 taps), so the thread that owns $\sigma=4$ carries 39\% of the
convolution work, and running scales concurrently leaves threads idle once the short scales finish. The
variants are equal up to $p=2$ and fall behind only at $p\ge4$.

\textbf{H4 (accuracy below lightweight deep models): supported.} The DRIVE test Dice of 0.7126 is below 0.8070
for LightVesselNet~\cite{sobhan2026}, 0.8328 for LFA-Net~\cite{mehmood2025lfa} and 0.8644 for
LVS-Net~\cite{mehmood2024lvs}. The pipeline needs no annotated data, no weights and no deep-learning runtime.
Its parameters transfer to STARE without retuning, losing 0.059 Dice, a drop driven mainly by images with
bright lesions. The deep-model figures are taken from the authors' papers and protocols, so the comparison is
indicative only.

\subsection{Why the small stages scale badly}
Morphology, thresholding and green-channel extraction each take well under a millisecond per image. A parallel
region on this machine costs a few microseconds to fork and join, which is a large fraction of a 0.2\,ms stage
split four ways, and at $p=8$ each logical CPU time-slices two threads, so a thread that is descheduled
delays the whole region. Morphology runs four such regions back to back, each over only 0.33\,MB of byte
data, so it slows down at $p=8$ (speedup 0.59). Because these stages make up under
0.5\% of the time, their poor scaling has no visible effect on the total. Fusing them into one parallel region
would be the fix if they ever mattered.

\subsection{Practical implications}
At 19.2 images per second on a 2-core VM, one such machine processes 10{,}000 fundus images in under nine
minutes, so throughput is not the limiting factor for a screening camp at DRIVE resolution. At 21 megapixels it
handles 0.26 images per second, where a second machine, more cores or a GPU would matter. Three practical
lessons follow. First, parallelise the Gaussian-derivative stage first, since it is 72\% of the time. Second,
on small SMT machines, check thread placement before adding threads: close binding with fewer threads than
logical CPUs can leave a whole core idle, while the default unbound run happened to spread correctly. Third,
prefer owner-computes histograms to atomics.

\subsection{Limitations}
The study uses one cloud VM with two physical cores shared with other tenants, so it says nothing about larger
core counts, NUMA systems or dedicated hardware, and its timings carry more noise than a dedicated machine's;
the coefficient of variation of the serial total was 2.6\%. The CPU model is reported only as an Intel Xeon at
2.20\,GHz (family 6, model 79). The placement and ablation experiments use five of the twenty test images. GPU
implementations, which the original proposal planned, are outside the scope of this paper. The accuracy
pipeline uses one global threshold and a fixed FOV erosion, and the STARE FOV masks are computed by
thresholding rather than supplied with the dataset.
